% GENERATED FILE. Edit canonical lessons, then run npm run generate:books.
% canonical-source-hash: fnv1a64:8b920c312b0b3cd0
% canonical-lessons: MW-C251-gardan, MW-C251-kadho, MW-C251-chhati, MW-C251-kohni, MW-C251-kamar

\chapter{Neck, Chest and Shoulders}
\label{ch:mw-neck-chest-and-shoulders}
\hlchaptermodality{251}

\begin{chapteropening}
\textbf{By the end of this chapter:} \emph{I can name more parts of the body.}

Complete the last lesson of chapter 251: I can name more parts of the body.
\end{chapteropening}

\section[\texorpdfstring{\mw{गरदन} (gardan)}{गरदन (gardan)}]{\mw{गरदन} (gardan) — the neck}

\label{lesson:MW-C251-gardan}



\begin{quote}
Before the new one: say the Marwadi for the waving of a lamp in worship, then the Marwadi for an uproar.
\end{quote}

\subsection*{You'll want to know: \mw{गरदन}}
\textbf{\mw{गरदन}} — \emph{gardan} — \textquotedblleft{}the neck\textquotedblright{}.

\begin{quote}
\mw{म्हारी} \mw{गरदन} \mw{लांबी} \mw{है।}
\end{quote}

\begin{quote}
\emph{mhārī gardan lā̃bī hai.}
\end{quote}

\begin{quote}
My neck is long.
\end{quote}

A feminine noun.

\begin{cousinweb}
\textbf{\mw{गरदन}} (\emph{gardan}) is usually traced to Persian \emph{gardan}, \textquotedblleft{}neck\textquotedblright{}.
\end{cousinweb}

\begin{grammarlens}[title={What you've built}]
One more word for this chapter.
\end{grammarlens}

\subsection*{Your turn}
\begin{itemize}
\raggedright
  \item \emph{Say it:} \emph{gardan}
  \item \emph{Say it:} \emph{gardan}, once more
  \item \emph{Say it:} \emph{mhārī gardan lā̃bī hai.} — My neck is long.
  \item \emph{Recall:} say the Marwadi for the waving of a lamp in worship, then the Marwadi for an uproar, then say \emph{gardan} again
\end{itemize}

\begin{tcolorbox}[breakable,colback=teal!4,colframe=teal!35!black,title={Before you move on}]
What does \mw{गरदन} mean? (\textquotedblleft{}The neck\textquotedblright{}.) Say it once more.
\end{tcolorbox}
\section[\texorpdfstring{\mw{कांधो} (kā̃dho)}{कांधो (kā̃dho)}]{\mw{कांधो} (kā̃dho) — a shoulder}

\label{lesson:MW-C251-kadho}



\begin{quote}
Before the new one: say the Marwadi for an uproar, then the Marwadi for the neck.
\end{quote}

\subsection*{You'll want to know: \mw{कांधो}}
\textbf{\mw{कांधो}} — \emph{kā̃dho} — \textquotedblleft{}a shoulder\textquotedblright{}.

\begin{quote}
\mw{म्हारो} \mw{कांधो} \mw{भारी} \mw{लागै} \mw{है।}
\end{quote}

\begin{quote}
\emph{mhāro kā̃dho bhārī lāgai hai.}
\end{quote}

\begin{quote}
My shoulder feels heavy.
\end{quote}

A masculine noun.

\begin{cousinweb}
\textbf{\mw{कांधो}} (\emph{kā̃dho}) is inherited from Sanskrit \emph{skandha}, \textquotedblleft{}shoulder\textquotedblright{}. CDIAL gives no Marwari form in this entry; it lists Hindi \emph{kā̃dhā} (CDIAL 13627).
\end{cousinweb}

\begin{grammarlens}[title={What you've built}]
One more word for this chapter.
\end{grammarlens}

\subsection*{Your turn}
\begin{itemize}
\raggedright
  \item \emph{Say it:} \emph{kā̃dho}
  \item \emph{Say it:} \emph{kā̃dho}, once more
  \item \emph{Say it:} \emph{mhāro kā̃dho bhārī lāgai hai.} — My shoulder feels heavy.
  \item \emph{Recall:} say the Marwadi for an uproar, then the Marwadi for the neck, then say \emph{kā̃dho} again
\end{itemize}

\begin{tcolorbox}[breakable,colback=teal!4,colframe=teal!35!black,title={Before you move on}]
What does \mw{कांधो} mean? (\textquotedblleft{}A shoulder\textquotedblright{}.) Say it once more.
\end{tcolorbox}
\section[\texorpdfstring{\mw{छाती} (chhātī)}{छाती (chhātī)}]{\mw{छाती} (chhātī) — the chest}

\label{lesson:MW-C251-chhati}



\begin{quote}
Before the new one: say the Marwadi for the neck, then the Marwadi for a shoulder.
\end{quote}

\subsection*{You'll want to know: \mw{छाती}}
\textbf{\mw{छाती}} — \emph{chhātī} — \textquotedblleft{}the chest\textquotedblright{}.

\begin{quote}
\mw{म्हारी} \mw{छाती} \mw{मांय} \mw{दरद} \mw{है।}
\end{quote}

\begin{quote}
\emph{mhārī chhātī mā̃y darad hai.}
\end{quote}

\begin{quote}
I have a pain in my chest.
\end{quote}

A feminine noun.

\begin{cousinweb}
\textbf{\mw{छाती}} (\emph{chhātī}) is inherited from a reconstructed \emph{*chātti}, \textquotedblleft{}chest, breast\textquotedblright{}; CDIAL lists Marwari \emph{chātī} under it (CDIAL 5014).
\end{cousinweb}

\begin{grammarlens}[title={What you've built}]
One more word for this chapter.
\end{grammarlens}

\subsection*{Your turn}
\begin{itemize}
\raggedright
  \item \emph{Say it:} \emph{chhātī}
  \item \emph{Say it:} \emph{chhātī}, once more
  \item \emph{Say it:} \emph{mhārī chhātī mā̃y darad hai.} — I have a pain in my chest.
  \item \emph{Recall:} say the Marwadi for the neck, then the Marwadi for a shoulder, then say \emph{chhātī} again
\end{itemize}

\begin{tcolorbox}[breakable,colback=teal!4,colframe=teal!35!black,title={Before you move on}]
What does \mw{छाती} mean? (\textquotedblleft{}The chest\textquotedblright{}.) Say it once more.
\end{tcolorbox}
\section[\texorpdfstring{\mw{कोहणी} (kohṇī)}{कोहणी (kohṇī)}]{\mw{कोहणी} (kohṇī) — an elbow}

\label{lesson:MW-C251-kohni}



\begin{quote}
Before the new one: say the Marwadi for a shoulder, then the Marwadi for the chest.
\end{quote}

\subsection*{You'll want to know: \mw{कोहणी}}
\textbf{\mw{कोहणी}} — \emph{kohṇī} — \textquotedblleft{}an elbow\textquotedblright{}.

\begin{quote}
\mw{म्हारी} \mw{कोहणी} \mw{माथै} \mw{घाव} \mw{है।}
\end{quote}

\begin{quote}
\emph{mhārī kohṇī māthai ghāv hai.}
\end{quote}

\begin{quote}
There is a wound on my elbow.
\end{quote}

A feminine noun.

\begin{cousinweb}
\textbf{\mw{कोहणी}} (\emph{kohṇī}) is inherited from Sanskrit \emph{kaphoṇi}, \textquotedblleft{}elbow\textquotedblright{}. CDIAL gives no Marwari form in this entry; it lists Hindi \emph{kohnī} (CDIAL 2757).
\end{cousinweb}

\begin{grammarlens}[title={What you've built}]
One more word for this chapter.
\end{grammarlens}

\subsection*{Your turn}
\begin{itemize}
\raggedright
  \item \emph{Say it:} \emph{kohṇī}
  \item \emph{Say it:} \emph{kohṇī}, once more
  \item \emph{Say it:} \emph{mhārī kohṇī māthai ghāv hai.} — There is a wound on my elbow.
  \item \emph{Recall:} say the Marwadi for a shoulder, then the Marwadi for the chest, then say \emph{kohṇī} again
\end{itemize}

\begin{tcolorbox}[breakable,colback=teal!4,colframe=teal!35!black,title={Before you move on}]
What does \mw{कोहणी} mean? (\textquotedblleft{}An elbow\textquotedblright{}.) Say it once more.
\end{tcolorbox}
\section[\texorpdfstring{\mw{कमर} (kamar)}{कमर (kamar)}]{\mw{कमर} (kamar) — the waist; the lower back}

\label{lesson:MW-C251-kamar}



\begin{quote}
Before the new one: say the Marwadi for the chest, then the Marwadi for an elbow.
\end{quote}

\subsection*{You'll want to know: \mw{कमर}}
\textbf{\mw{कमर}} — \emph{kamar} — \textquotedblleft{}the waist; the lower back\textquotedblright{}.

\begin{quote}
\mw{म्हारी} \mw{कमर} \mw{पतळी} \mw{है।}
\end{quote}

\begin{quote}
\emph{mhārī kamar patḷī hai.}
\end{quote}

\begin{quote}
My waist is slim.
\end{quote}

A feminine noun.

\begin{cousinweb}
\textbf{\mw{कमर}} (\emph{kamar}) is usually traced to Persian \emph{kamar}, \textquotedblleft{}waist\textquotedblright{}.
\end{cousinweb}

\begin{grammarlens}[title={What you've built}]
One more word for this chapter.
\end{grammarlens}

\subsection*{Your turn}
\begin{itemize}
\raggedright
  \item \emph{Say it:} \emph{kamar}
  \item \emph{Say it:} \emph{kamar}, once more
  \item \emph{Say it:} \emph{mhārī kamar patḷī hai.} — My waist is slim.
  \item \emph{Recall:} say the Marwadi for the chest, then the Marwadi for an elbow, then say \emph{kamar} again
\end{itemize}

\begin{tcolorbox}[breakable,colback=teal!4,colframe=teal!35!black,title={Before you move on}]
What does \mw{कमर} mean? (\textquotedblleft{}The waist; the lower back\textquotedblright{}.) Say it once more.
\end{tcolorbox}
